\section*{अध्याय \ref{ch:b2:biogeochemical-cycles} --- जैव-भू-रासायनिक चक्र}

\begin{solution}{exo:b2:biogeochemical-cycles:1}
कोई \omterm{def:b2:biogeochemical-cycles:reservoir}{भंडार} उस तत्त्व का कोई संचय (कोई द्रव्यमान) है; कोई \omterm{def:b2:biogeochemical-cycles:reservoir}{अभिवाह}
\omterm{def:b2:biogeochemical-cycles:reservoir}{भंडारों} के बीच संचरण की कोई दर (प्रति वर्ष द्रव्यमान); और \omterm{def:b2:biogeochemical-cycles:reservoir}{निवास-काल}
उस \omterm{def:b2:biogeochemical-cycles:reservoir}{भंडार} का द्रव्यमान बँटा उसका कुल बहिर्वाह है, यानी वह माध्य समय जो कोई
परमाणु वहाँ बिताता है। स्थल-पौधे: $450/120 \approx 4$ वर्ष (मोटे तौर पर किसी पत्ती का
कार्बन; काठ दशकों रहता है और वह औसत उस
फैलाव को छिपा देता है)। गहरा महासागर: $37000/100 \approx 400$ वर्ष --- यानी वह समय
जो किसी डूबे अणु को सतह फिर देखने में लगता है।
\end{solution}

\begin{solution}{exo:b2:biogeochemical-cycles:2}
वर्ष भर में $2.5\times 2.13 = \qty{5.3}{GtC}$। $10/2.13 = \qty{4.7}{ppm}$
--- यदि वह सब वायु में रहे; और \qty{45}{\%} के \omterm{prop:b2:biogeochemical-cycles:carbon}{वायुवाहित
अंश} पर, \qty{2.1}{ppm}।
\end{solution}

\begin{solution}{exo:b2:biogeochemical-cycles:3}
खोलना: नाइट्रोजनेज़ से जैविक स्थिरीकरण (सायनोजीवाणु,
राइजोबिया, \emph{Azotobacter}, \emph{Frankia}) और बिजली। वापसी:
\omterm{def:b2:microbial-metabolism:anaerobic}{विनाइट्रीकरण} (अनॉक्सी मिट्टी तथा अवसाद में \emph{Pseudomonas}
और \emph{Paracoccus} जैसे विकल्पी अवायुजीवी) और ऐनैमॉक्स
(प्लैंक्टोमाइसीट)। बीच में: अपघटकों (कवक, जीवाणु) द्वारा कार्बनिक नाइट्रोजन का
अमोनियम में अमोनीकरण; अमोनियम का
नाइट्राइट में नाइट्रीकरण (\emph{Nitrosomonas}, अमोनिया-ऑक्सीकारी
आर्किया) तथा नाइट्राइट का नाइट्रेट में (\emph{Nitrobacter}); और
पौधों तथा सूक्ष्मजीवों द्वारा अमोनियम या नाइट्रेट का स्वांगीकरण।
\end{solution}

\begin{solution}{exo:b2:biogeochemical-cycles:4}
भूवैज्ञानिक समय पर नाइट्रोजन सदा वायु की
अक्षय $\mathrm{N_2}$ से स्थिरीकारियों द्वारा बनाई जा सकती है, जो जब भी नाइट्रोजन कम हो तब
अनुकूल पाए जाते हैं; जबकि फ़ॉस्फ़ोरस का ऐसा कोई \omterm{def:b2:biogeochemical-cycles:reservoir}{भंडार} नहीं है और
उसकी आपूर्ति अपक्षय से नियत है, इसलिए उत्पादकता पर दीर्घकालिक छत
फ़ॉस्फ़ोरस है। किसी ऋतु के भीतर स्थिरीकरण धीमा तथा
महँगा है, और अधिकांश पौधे वह कर ही नहीं सकते, इसलिए जो नाइट्रोजन हाथ में हो वही
यहाँ तथा अभी वृद्धि सीमित करती है --- और इसीलिए उर्वरक
अधिकतर नाइट्रोजन है।
\end{solution}

\begin{solution}{exo:b2:biogeochemical-cycles:5}
$\tau = 1/0.2 = 5$ वर्ष; $M^{*} = F\tau = \qty{10}{t}$, यानी
$\qty{1}{g/m^3} = \qty{1000}{mg/m^3}$, यानी घोर सुपोषी।
\qty{90}{\%} तक का काल: $\tau\ln 10 = 11.5$ वर्ष। उस निवेश को आधा करना
उस \omterm{def:b2:biogeochemical-cycles:reservoir}{स्थायी अवस्था} को आधा कर देता है, यानी \qty{500}{mg/m^3}, और वह भी उसी
काल-मान पर।
\end{solution}

\begin{solution}{exo:b2:biogeochemical-cycles:6}
1960 का दशक: \qty{0.9}{ppm/yr}, \qty{1.9}{GtC/yr}। 2010 का दशक: \qty{2.3}{ppm/yr},
\qty{5.0}{GtC/yr}। वह वृद्धि-दर उत्सर्जनों के साथ क़दम मिलाकर ढाई
गुना बढ़ चुकी है।
\end{solution}

\begin{solution}{exo:b2:biogeochemical-cycles:7}
\qty{6}{ppm} यानी \qty{13}{GtC}, जबकि शुद्ध प्राथमिक उत्पादन
\qty{35}{GtC} है: यानी वह आरा-दाँत उसका एक तिहाई। वह इसलिए छोटा है कि
श्वसन तथा अपघटन ग्रीष्म भर चलते रहते हैं (वह
आरा-दाँत शुद्ध ग्रहण अभिलिखित करता है, सकल नहीं), क्योंकि
दक्षिणी गोलार्ध की ऋतुएँ उलटी हैं और आंशिक रूप से कट जाती हैं, और
क्योंकि वह महासागर उस झूल को कुंद करता है।
\end{solution}

\begin{solution}{exo:b2:biogeochemical-cycles:8}
$\qty{200}{kg}/(\qty{28}{g/mol}) = 7100$~मोल $\mathrm{N_2}$; $16\times
7100 = \num{114000}$~मोल ATP; $/30 = 3800$~मोल ग्लूकोज़, \qty{690}{kg}।
जो फ़सल \qty{10}{t} शुष्क पदार्थ बनाती है वह श्वसन गिनने पर मोटे तौर पर
\qtyrange{15}{20}{t} शर्करा स्थिर करती है, इसलिए
स्थिरीकरण का दाम उसके प्रकाशसंश्लिष्ट का कोई \qty{4}{\%} है --- यानी
मिट्टी की नाइट्रोजन से स्वतंत्रता का दाम।
\end{solution}

\begin{solution}{exo:b2:biogeochemical-cycles:9}
उपलब्ध अनुपात $\mathrm{N:P} = 32/2.2 = 14.5$, जो रेडफ़ील्ड के
16 से नीचे है: नाइट्रेट पहले चुकता है, और \qty{0.2}{\micro mol/kg}
फ़ॉस्फ़ेट बचा रह जाता है। स्थिर कार्बन: $32\times 106/16 =
\qty{212}{\micro mol/kg}$, यानी प्रति किलोग्राम जल लगभग \qty{2.5}{mg}
कार्बन।
\end{solution}

\begin{solution}{exo:b2:biogeochemical-cycles:10}
$E^{*} = Q\tau = \qty{500}{GtC}$, यानी उद्योग-पूर्व 285 से
\qty{235}{ppm} ऊपर: यानी लगभग \qty{520}{ppm}। असल में उससे बुरा, क्योंकि
वे पात्र कोई एक ही तेज़ कोष नहीं हैं: वह सतही महासागर संतृप्त होता है
(उसकी $\mathrm{CO_2}$ लेने की क्षमता उसके कार्बोनेट के चुकने के साथ गिरती है),
वह स्थल-पात्र तापन के साथ कमज़ोर पड़ता है, और उस सच्चे हटाव का कोई
धीमा घटक है --- गहरे-महासागर का मिश्रण तथा चट्टान-अपक्षय --- जिसके
काल-नियतांक हज़ार से एक लाख वर्ष के हैं, इसलिए उस अतिरेक का कोई
अंश दशकों के किसी भी $\tau$ पर शिथिल नहीं होता।
\end{solution}

\begin{solution}{exo:b2:biogeochemical-cycles:11}
जो मापा गया वह वायुमंडलीय $\mathrm{CO_2}$ का महासागर तथा जीवमंडल से
विनिमय है, क्षय नहीं: वह बम-$\mathrm{{}^{14}C}$
सतही महासागर तथा पौधों और मिट्टियों के कहीं बड़े कोषों में तनु कर दिया गया,
जिन्होंने फिर साधारण कार्बन लौटाया। उस अतिरेक की 11-वर्ष अर्ध-आयु
(यानी लगभग 16 वर्ष का कोई e-गुना काल) उस विनिमय का
काल-मान है --- यानी सतही महासागर तथा स्थल के सापेक्ष वायुमंडल के कार्बन का
फेर, जो उस चार-वर्ष सकल \omterm{def:b2:biogeochemical-cycles:reservoir}{निवास-काल} से लंबा है क्योंकि जो निकलता है उसका
अधिकांश सीधे लौट आता है।
\end{solution}

\begin{solution}{exo:b2:biogeochemical-cycles:12}
नाइट्रोजन: प्राकृतिक स्थिरीकरण वर्ष भर में लगभग 200~Tg; मानवीय स्थिरीकरण
(हाबर--बॉश 120, फ़सलें 60, दहन 30) लगभग 210: यानी वह \omterm{def:b2:biogeochemical-cycles:reservoir}{अभिवाह}
दुगुना हो चुका है, और प्रभावित \omterm{def:b2:biogeochemical-cycles:reservoir}{भंडार} --- मिट्टियाँ, नदियाँ, तटीय
समुद्र, वायु की $\mathrm{N_2O}$ --- छोटे तथा तेज़ हैं, इसलिए वह परिवर्तन
दशकों के भीतर दिखता है, जबकि वह $\mathrm{N_2}$ \omterm{def:b2:biogeochemical-cycles:reservoir}{भंडार} अछूता है।
कार्बन: वर्ष भर में 11~GtC के मानवीय उत्सर्जन लगभग 200 के
सकल प्राकृतिक विनिमय का केवल \qty{5}{\%} हैं, इसलिए वह \omterm{def:b2:biogeochemical-cycles:reservoir}{अभिवाह} ठेला भर गया है;
पर वे प्राकृतिक \omterm{def:b2:biogeochemical-cycles:reservoir}{अभिवाह} कट जाते हैं और वह मानवीय नहीं कटता, और उसने
वायुमंडलीय \omterm{def:b2:biogeochemical-cycles:reservoir}{भंडार} \qty{50}{\%} चढ़ा दिया है --- यानी उस छोटे
ठेले का किसी ऐसे \omterm{def:b2:biogeochemical-cycles:reservoir}{भंडार} पर बड़ा संचयी प्रभाव है जिसका अंतिम पात्र धीमा
है। कौन-सा चक्र ``अधिक विक्षुब्ध'' है यह इस पर निर्भर है कि
कोई \omterm{def:b2:biogeochemical-cycles:reservoir}{अभिवाह} गिने या संचय।
\end{solution}

\begin{solution}{pb:b2:biogeochemical-cycles:1}
\textbf{1.} $285\times 2.13 = \qty{607}{GtC}$; $412\times 2.13 =
\qty{878}{GtC}$; वृद्धि \qty{271}{GtC}।
\textbf{2.} $271/700 = 0.39$ (हाल के दशकों भर वह अंश,
भू-उपयोग उत्सर्जन गिनने पर, $0.45$ के अधिक पास है)।
\textbf{3.} वह महासागर, यानी उत्सर्जनों का लगभग एक चौथाई, और वह स्थल
(वन-पुनर्वृद्धि तथा $\mathrm{CO_2}$ उर्वरण), लगभग एक तिहाई; यानी उस
\qty{430}{GtC} का शेष जो वायु में नहीं है।
\textbf{4.} $175/38000 = \qty{0.5}{\%}$ --- यानी पूरे महासागर के लिए
नगण्य; पर वह ग्रहण अब तक मुख्यतः सतही परत
(\qty{900}{GtC}) में घुसा है, जहाँ वह दसवें भाग या उससे अधिक की चढ़ाई है, और
वह कार्बोनेट रसायन घुली $\mathrm{CO_2}$ में किसी छोटी चढ़ाई को
हाइड्रोजन-आयन सांद्रता में किसी बड़ी चढ़ाई में बदल देता है: सतही pH
$0.1$ गिर चुका है, यानी अम्लता में \qty{30}{\%} की चढ़ाई।
\textbf{5.} $Q = Q_0 e^{t/T}$ के साथ, $E = Ae^{t/T}$ आज़माइए: $A/T = Q_0 -
A/\tau$, इसलिए $A = Q_0\tau T/(\tau + T)$ और वह \omterm{prop:b2:biogeochemical-cycles:carbon}{वायुवाहित अंश}
$\dot E/Q = \tau/(\tau + T)$ अचल है। वर्ष भर में
\qty{2}{\%} से बढ़ते उत्सर्जन ($T = 50$) $\tau = 40$ के साथ $0.44$ देते हैं: वह
अंश इसलिए अचल है कि पात्र तथा स्रोत साथ बढ़ते हैं।
\textbf{6.} $Q$ अचल हो, तो $E \to Q\tau$ और $\dot E \to 0$: यानी वह
\omterm{prop:b2:biogeochemical-cycles:carbon}{वायुवाहित अंश} शून्य की ओर गिरता है ज्यों-ज्यों वे पात्र पकड़ते जाएँ ---
उस एक-डिब्बा प्रतिरूप में; असल में वे पात्र संतृप्त होते हैं और वह अंश
अधिक धीरे गिरता है।
\textbf{7.} \qty{271}{GtC}।
\textbf{8.} $E^{*} = 10\times 100 = \qty{1000}{GtC}$: $285 +
1000/2.13 = \qty{755}{ppm}$।
\textbf{9.} $E(t) = 1000 + (271 - 1000)e^{-t/100}$; $t = 80$ पर,
$E = 1000 - 729\times 0.449 = \qty{672}{GtC}$: \qty{600}{ppm}।
\textbf{10.} $Q = 10(1 - t/50)$: विशेष हल $E_p = at + b$
जिसमें $a = -20$, $b = 3000$ ($a = 10 - b/\tau$ तथा $0 = -0.2 -
a/\tau$ से); $E = 3000 - 20t + (271 - 3000)e^{-t/100}$। $t = 50$ पर:
$2000 - 2729\times 0.607 = \qty{344}{GtC}$, \qty{447}{ppm}। फिर $E$
स्वतंत्र रूप से क्षीण होता है: $t = 80$ पर, $344\,e^{-0.3} = \qty{255}{GtC}$,
\qty{405}{ppm}।
\textbf{11.} \qty{405}{ppm} के मुक़ाबले \qty{600}{ppm}: यानी दूसरा
परिदृश्य उस वायु को, 2100 तक, लगभग वहीं लौटा देता है जहाँ वह आज है।
\textbf{12.} वह एक ही $\tau$ उन पात्रों से किसी सिकुड़ते अतिरेक पर
उसी दर से काम करवाता रहता है; असल में वे तेज़ पात्र (वह
मिश्रित परत, वे बढ़ते वन) जो ले सकते थे वह पहले ही ले चुके हैं
और बाक़ी हटाव शताब्दियों से सहस्राब्दियों के $\tau$ वाले धीमे गहरे-महासागर मिश्रण तथा
अपक्षय से होता है, इसलिए शुद्ध शून्य के बाद वह ह्रास कहीं धीमा है ---
यानी वह सांद्रता गिरने के बजाय शताब्दियों तक मोटे तौर पर
अचल रहती है।
\textbf{13.} फ़सल \qty{90}{kg}; विनाइट्रीकृत \qty{37.5}{kg}; बही
\qty{22.5}{kg} नाइट्रोजन प्रति हेक्टेयर प्रति वर्ष।
\textbf{14.} $\qty{22.5}{kg}/\qty{14}{g/mol} = 1600$~मोल N; कार्बन
$1600\times 106/16 = \num{10600}$~मोल, \qty{128}{kg}।
\textbf{15.} \num{10600}~मोल $\mathrm{O_2}$, \qty{340}{kg}।
\textbf{16.} $\num{10600}/0.25 = \num{42000}$~\unit{m^3}: यानी एक
हेक्टेयर का बहाव हर वर्ष तली के चालीस हज़ार घन
मीटर जल से ऑक्सीजन छीन सकता है --- और मिसिसिपी
दस करोड़ हेक्टेयर बहाती है।
\textbf{17.} $37.5\times 0.01 = \qty{0.375}{kg}$ नाइट्रोजन
$\mathrm{N_2O}$ के रूप में, यानी $0.375\times 44/28 = \qty{0.59}{kg}$ $\mathrm{N_2O}$;
$\times 270 = \qty{160}{kg}$ $\mathrm{CO_2}$ तुल्यांक।
\textbf{18.} $150\times 4 = \qty{600}{kg}$ $\mathrm{CO_2}$: यानी वह
निर्माण उस खेत की $\mathrm{N_2O}$ से चार गुना भारी है, और
दोनों मिलकर प्रति हेक्टेयर प्रति वर्ष पौन टन हैं।
\textbf{19.} $\tau = 10$ वर्ष; $M^{*} = 5\times 10 = \qty{50}{t}$;
सांद्रता $50/(5\times 10^{7}) = \qty{1}{g/m^3} =
\qty{1000}{mg/m^3}$।
\textbf{20.} उस देहली से तीस गुना ऊपर: यानी घोर सुपोषी।
\textbf{21.} $1/\tau = 0.1 + 0.2 = 0.3$: $\tau = 3.3$ वर्ष; $M^{*} =
5/0.3 = \qty{16.7}{t}$, \qty{330}{mg/m^3}।
\textbf{22.} $M^{*} = 1/0.3 = \qty{3.3}{t}$, \qty{67}{mg/m^3} --- यानी अब भी
सुपोषी; और $\tau\ln 10 = 7.7$ वर्ष बाद \qty{10}{\%} के भीतर।
\textbf{23.} वह अवसाद कोई एकतरफ़ा पात्र नहीं है: उन सुपोषी वर्षों में
उसमें संचित फ़ॉस्फ़ोरस उन्हीं अनॉक्सी दशाओं के नीचे वापस छोड़ा जाता है
जो सुपोषण स्वयं बनाता है (लोहे से बँधा
फ़ॉस्फ़ेट लोहे के अपचयित होने पर घुल जाता है), इसलिए वह झील उस बाहरी निवेश के
कटने के बाद भी वर्षों तक स्वयं को पोसती है --- यानी भीतरी
भरण, कोई दूसरा \omterm{def:b2:biogeochemical-cycles:reservoir}{भंडार} लंबे \omterm{def:b2:biogeochemical-cycles:reservoir}{निवास-काल} वाला।
\textbf{24.} $\qty{5000}{kg}/\qty{31}{g/mol} = 1.6\times 10^{5}$~मोल P;
कार्बन $\times 106 = 1.7\times 10^{7}$~मोल, यानी वर्ष भर में \qty{205}{t}
कार्बन --- यानी उस झील की सतह पर \qty{20}{g/m^2} यदि वह
\qty{5}{m} गहरी हो, यानी कोई भारी प्रस्फुटन।
\textbf{25.} 1850--2020 भर \omterm{prop:b2:biogeochemical-cycles:carbon}{वायुवाहित अंश} $0.39$; 2100 में,
\qty{10}{GtC/yr} पर थामे गए उत्सर्जनों के साथ \qty{600}{ppm} और
2070 तक शून्य तक काटे गए उत्सर्जनों के साथ \qty{405}{ppm} (यानी कोई आशावादी
एकल-$\tau$ प्रतिरूप); और बही हुई नाइट्रोजन प्रति हेक्टेयर प्रति वर्ष
\qty{22.5}{kg}।
\end{solution}
